% Two install disciplines in, one resolve path out.
%
% This figure exists to DRAW prop:kindinvisible rather than argue it: every
% difference between the kinds is on the ARROW from slot to record, and nothing
% past the record can tell which one parked a slot.  If a reader takes only that
% away, the figure has done its job.
%
% Same visual vocabulary as fig:tristate and P1's fig:fliplatch -- deliberately.
%
% The slots are drawn IDENTICALLY on purpose.  It is tempting to colour the
% exclusive one differently, but that would be untrue to the mechanism: a slot's
% kind is not a property the word carries, it is a property of the WRITER's
% discipline, so it belongs on the arrow and nowhere else.
%
% Deliberately NOT drawn: the install ORDER (shared records first, then the
% exclusive parks, then the status store, then settle).  That is sequencing, it
% belongs in prose, and crowding it in here would bury the single point above.
%
% Layout note: an earlier draft ran a rotated "the kinds differ here" label down
% the middle and it overprinted the arrow labels.  The two shared slots sit close
% together and the exclusive one sits apart, which separates the kinds by
% POSITION instead, leaving the top clear for the failure edge.  Render the page
% and look at it before adding anything.
\begin{figure}[t]
\centering
\begin{tikzpicture}[
  font=\small,
  slot/.style   = {draw, minimum width=15mm, minimum height=6mm, fill=gray!8},
  proxy/.style  = {draw, minimum width=26mm, minimum height=9mm, fill=blue!6},
  target/.style = {draw, rounded corners=2pt, minimum width=11mm, minimum height=5.5mm},
  old/.style    = {target, fill=gray!15},
  sel/.style    = {draw, very thick, minimum width=17mm, minimum height=8mm, fill=orange!18},
  arr/.style    = {-latex, thin},
  dep/.style    = {-latex, thin, densely dotted, gray!70!black},
  fail/.style   = {-latex, thin, densely dashed, red!60!black},
]

% ---- slots: drawn identically; the kind is not visible in the slot ----------
\node[slot] (s1) at (0,  1.7) {\code{s$_1$}};
\node[slot] (s2) at (0,  0.5) {\code{s$_2$}};
\node[slot] (s3) at (0, -1.5) {\code{s$_3$}};
\node[left=1.5mm of s1, font=\scriptsize\itshape] {shared};
\node[left=1.5mm of s2, font=\scriptsize\itshape] {shared};
\node[left=1.5mm of s3, font=\scriptsize\itshape, align=right] {exclusive\\(lock held)};

% ---- records ---------------------------------------------------------------
\node[proxy] (p1) at (4.4,  1.7) {$\{o_1,\, n_1\}$};
\node[proxy] (p2) at (4.4,  0.5) {$\{o_2,\, n_2\}$};
\node[proxy] (p3) at (4.4, -1.5) {$\{o_3,\, n_3\}$};

% ---- the two installs: the ONLY place the kinds differ ---------------------
\draw[arr] (s1) -- node[above, font=\scriptsize] {CAS} (p1);
\draw[arr] (s2) -- node[above, font=\scriptsize] {CAS} (p2);
\draw[arr, thick] (s3) --
  node[above, font=\scriptsize] {plain store}
  node[below, font=\scriptsize] {cannot fail} (p3);

% The failure edge belongs to the CAS installs and to nothing else.  Anchored to
% the arrow's own midpoint, drawn up into space kept clear for it.
\draw[fail] (2.2, 1.95) -- ++(0.5, 0.75)
  node[right, font=\scriptsize, text=red!60!black] {a peer changed it: \emph{abort}};

% ---- one status word -------------------------------------------------------
\node[sel] (st) at (8.8, 0.2) {\code{status}};
\draw[dep] (p1.east) -- (st.north west);
\draw[dep] (p2.east) -- (st.west);
\draw[dep] (p3.east) -- (st.south west);

% ---- one resolve path, identical for both kinds ----------------------------
\node[old] (res) at (12.6, 0.2) {$o_i$ / $n_i$};
\draw[arr, thick] (st.east) -- node[above, font=\scriptsize] {resolve} (res.west);

% Short, and tucked directly under the resolve arrow it labels.  Two earlier
% placements collided with the dotted arrows converging on `status` from the
% records -- that whole wedge of space is crossed by them and is not usable for
% text.  Under the resolve arrow, right of `status`, nothing crosses.  The long
% form of this sentence lives in the caption, where it cannot collide with
% anything.
\node[font=\footnotesize\itshape, text=gray!45!black]
  at (10.85, -0.5) {identical for both kinds};

\end{tikzpicture}
\caption{Two install disciplines, one resolve path. A slot with concurrent
writers is planted with a compare-and-swap, which a peer's intervening change
makes fail---the only source of a contention abort. A slot the embedder holds a
lock over is parked with a plain release store, which cannot fail. The slots are
drawn identically because they \emph{are} identical: the kind is a property of
the writer's discipline, not something the word carries, so it lives on the
arrow and nowhere else. Past the records the two are one object, and a reader
resolving a slot cannot determine which discipline parked it. A transaction with
no shared record therefore has nothing in it that can abort.}
\label{fig:kinds}
\end{figure}
